\hypertarget{how-to-manage-non-personal-system-accounts}{%
\section{How to manage non personal (system)
accounts}\label{how-to-manage-non-personal-system-accounts}}

\begin{wrapfigure}{i}{0.3\textwidth}
\centering
\includegraphics{./10_accounts/10_non_personal_accounts/b4974b37d04e9c5e7e898d82c65a9566fdf029fe.jpg}
\caption{André Koot}
\end{wrapfigure}

Customer often ask me for best practices regarding management of
non-personal or highly privileged accounts in the process of
implementing an Identity and Access Management (IAM) solution. This is
an interesting question, because in an IAM project, we try to manage all
kinds of accounts, but this type of account is different from accounts
that are owned by end users. This type of accounts can't directly be
related to a uniquely identifiable person, nor are they the result of
the `joiner -- mover - leaver' HR processes in an organization. So, how
do you manage the existence of such an account?

\textbf{Types of non personal accounts}

There are Non Personal Accounts (NPA's) and Non-personal System
Accounts: (NPSA's). We can identify:

- \textbf{Admin or root account}

The admin or root account of Windows and Linux or Unix servers is highly
privileged system account on the respective platforms.

o It is authorized at the highest level

o It has access to every file and process running on a platform.

o The `root' or `Admin' has the permissions to change the behavior of
their component;

o Commands can be run from it as well as react to responses of the
system.

o Operational use of the account needs to be monitored continuously.

- \textbf{Superuser account} It's a business information system or
application account, that looks a lot like `root'. It is there when the
system is installed, it's a system account. The Superuser has permission
to modify, making it a risk critical account in an information system.
Like Sap* in a Sap environment.

- \textbf{Service account}

Accounts for middleware processes like DBMS's, ESB's or other ICT
components that run on top of the Windows or Linux operating systems. A
special form of a non-personal account is an application account in a
DBMS to give database access to an application.

- \textbf{Batch user account}

An account used by a batch job process, it is most commonly used for
scheduled batch jobs, like nightly file transfers.

\textbf{NPA characteristics}

NPSA's have a few characteristics in common. They are non-personal and
they are not directly connected to a person. Login with the account
doesn't leave an audit trace showing which person is actually using it.
And, of course, NPSA's are very powerful and so use of them should be
tightly controlled.

Service and batch accounts also have a specific similarity: one
typically doesn't login with such an account, these accounts are not
used interactively. In most cases such an account is only used as a
placeholder with some permissions to perform specific tasks, like
running a webserver with the limited capability to process
https-requests and write log files.

Modern IAM solutions can be implemented to facilitate provisioning of
personal accounts for specific user functionality. Since non-personal
accounts are not an attribute of an identity, there is not a sole user
that can be connected to an NPA, so an IAM solution is not suitable to
manage them.

\textbf{If not an IAM problem, then what?}

These accounts belong to the component that they manage. The Windows
operating system comes with the Administrator account, Linux comes with
root. You cannot install Linux without a root account. You may not by
default be able to login with it (as on Ubuntu), but the account is
there. So by installing an OS, you automatically get the `God'‑account.
There is no choice, it is the result of the change process that leads to
the implementation of the OS. Such an NPSA should only be used in a
controlled manner from specific processes, like an incident management
process (admin may be needed to assist in a catastrophe), or the change
management process (the admin permissions may be required to perform an
infrastructural change).

The same is true for service accounts: when installing a middleware
component, like a database management system, the account is created to
enable the service, hence the name service account. Again, you have no
choice. You might install the service using a `root' -- type account,
but that will result in a security violation, Thou shalt not run any
service as root!

And again, for batch account the same is true again: a batch process is
created as the result of a change request. The batch tasks are created
to support an information system, or a business process. The batch job
is created to make it possible to schedule the automatic execution of
the tasks. A batch account is created to make it possible to use
resources on the system.

This leads to the following conclusion:

\textbf{Non-personal accounts have to be managed in the change
management process.}

\textbf{This has the following implications:}

\begin{itemize}
\tightlist
\item
  The account has to be registered in the configuration management
  database, it is an attribute of the component that it belongs to.
  Admin belongs to the Active Directory. Root belongs to a linux server.
  The account named `Oracle' probably belongs to an Oracle dbms
  instance: the dbms is a managed component and the account name is
  Oracle.
\item
  The account has an owner, who is accountable for the use of the
  account. Admin and root belong to the manager of the ICT department.
  The SAP account is owned by the system owner of the SAP system
\item
  The interactive accounts should only be used for infrastructural
  changes or calamities.
\item
  The non-personal system accounts should never be used interactively
  for operational tasks.
\item
  The passwords of these accounts must remain secret. They should be
  secured by means of an envelope procedure, a password vault, or by
  using a Privileged Account Management system (PAM, like CyberArk,
  Hitachi PAM or CA PAM to name just a few). Any use has to be related
  to a service management ticket (an incident or a change).
\end{itemize}

So, there you have it. Non-personal accounts must not be managed in an
IAM solution, they have to be managed by the Change Management processes
in an organization. Either they are owned by ICT or by the system owner
who owns the information system that the account is used for. You should
not manage privileged accounts in an IAM solution. And if you have to
execute tasks with one of these accounts: use a Privileged Account
Management system to secure it.
